% problem_id: S001-lemma-cech-Hausdorff-quasi-compact
% source_id: S001 (Stacks Project)
% source_file: cohomology.tex
% statement_locator: cohomology.tex lines 2666--2673
% proof_locator: cohomology.tex lines 2675--2845
% env: lemma
% id_note: label 在本来源内唯一
% pairing: tight (verified)
% status: complete (statement + author proof verbatim + reason + locators)
%
\begin{lemma}
\label{lemma-cech-Hausdorff-quasi-compact}
Let $X$ be a Hausdorff and quasi-compact topological space. Let
$\mathcal{F}$ be an abelian sheaf on $X$. Then
the map $\check{H}^n(X, \mathcal{F}) \to H^n(X, \mathcal{F})$ defined
in (\ref{equation-cech-to-cohomology}) is an isomorphism for
all $n$.
\end{lemma}

\begin{proof}
We already know that $\check{H}^n(X, -) \to H^n(X, -)$
is an isomorphism of functors for $n = 0, 1$, see
Lemma \ref{lemma-cech-always}.
The functors $H^n(X, -)$ form a universal $\delta$-functor, see
Derived Categories, Lemma \ref{derived-lemma-higher-derived-functors}.
If we show that $\check{H}^n(X, -)$ forms a universal $\delta$-functor
and that $\check{H}^n(X, -) \to H^n(X, -)$ is compatible with boundary
maps, then the map will automatically be an isomorphism by uniqueness
of universal $\delta$-functors, see
Homology, Lemma \ref{homology-lemma-uniqueness-universal-delta-functor}.

\medskip\noindent
Let $0 \to \mathcal{F} \to \mathcal{G} \to \mathcal{H} \to 0$
be a short exact sequence of abelian sheaves on $X$.
Let $\mathcal{U} : X = \bigcup_{i \in I} U_i$ be an open covering.
This gives a complex of complexes
$$
0 \to \check{\mathcal{C}}^\bullet(\mathcal{U}, \mathcal{F}) \to
\check{\mathcal{C}}^\bullet(\mathcal{U}, \mathcal{G}) \to
\check{\mathcal{C}}^\bullet(\mathcal{U}, \mathcal{H}) \to 0
$$
which is in general not exact on the right. The sequence defines
the maps
$$
\check{H}^n(\mathcal{U}, \mathcal{F}) \to
\check{H}^n(\mathcal{U}, \mathcal{G}) \to
\check{H}^n(\mathcal{U}, \mathcal{H})
$$
but isn't good enough to define a boundary operator
$\delta : \check{H}^n(\mathcal{U}, \mathcal{H}) \to
\check{H}^{n + 1}(\mathcal{U}, \mathcal{F})$. Indeed
such a thing will not exist in general. However, given an
element $\overline{h} \in \check{H}^n(\mathcal{U}, \mathcal{H})$
which is the cohomology class of a cocycle
$h = (h_{i_0 \ldots i_n})$
we can choose open coverings
$$
U_{i_0 \ldots i_n} = \bigcup W_{i_0 \ldots i_n, k}
$$
such that $h_{i_0 \ldots i_n}|_{W_{i_0 \ldots i_n, k}}$
lifts to a section of $\mathcal{G}$ over $W_{i_0 \ldots i_n, k}$.
By Topology, Lemma \ref{topology-lemma-refine-covering}
(this is where we use the assumption that $X$ is hausdorff and quasi-compact)
we can choose an open covering $\mathcal{V} : X = \bigcup_{j \in J} V_j$
and $\alpha : J \to I$ such that $V_j \subset U_{\alpha(j)}$
(it is a refinement) and such that for all $j_0, \ldots, j_n \in J$
there is a $k$ such that
$V_{j_0 \ldots j_n} \subset W_{\alpha(j_0) \ldots \alpha(j_n), k}$.
We obtain maps of complexes
$$
\xymatrix{
0 \ar[r] &
\check{\mathcal{C}}^\bullet(\mathcal{U}, \mathcal{F}) \ar[d] \ar[r] &
\check{\mathcal{C}}^\bullet(\mathcal{U}, \mathcal{G}) \ar[d] \ar[r] &
\check{\mathcal{C}}^\bullet(\mathcal{U}, \mathcal{H}) \ar[d] \ar[r] &
0 \\
0 \ar[r] &
\check{\mathcal{C}}^\bullet(\mathcal{V}, \mathcal{F}) \ar[r] &
\check{\mathcal{C}}^\bullet(\mathcal{V}, \mathcal{G}) \ar[r] &
\check{\mathcal{C}}^\bullet(\mathcal{V}, \mathcal{H}) \ar[r] &
0
}
$$
In fact, the vertical arrows are the maps of complexes used
to define the transition maps between the {\v C}ech cohomology groups.
Our choice of refinement shows that we may choose
$$
g_{j_0 \ldots j_n} \in
\mathcal{G}(V_{j_0 \ldots j_n}),\quad
g_{j_0 \ldots j_n} \longmapsto
h_{\alpha(j_0) \ldots \alpha(j_n)}|_{V_{j_0 \ldots j_n}}
$$
The cochain $g = (g_{j_0 \ldots j_n})$ is not a cocycle
in general but we know that its {\v C}ech boundary $\text{d}(g)$
maps to zero in $\check{\mathcal{C}}^{n + 1}(\mathcal{V}, \mathcal{H})$
(by the commutative diagram above and the fact that $h$ is a cocycle).
Hence $\text{d}(g)$ is a cocycle in
$\check{\mathcal{C}}^\bullet(\mathcal{V}, \mathcal{F})$.
This allows us to define
$$
\delta(\overline{h}) = \text{class of }\text{d}(g)\text{ in }
\check{H}^{n + 1}(\mathcal{V}, \mathcal{F})
$$
Now, given an element $\xi \in \check{H}^n(X, \mathcal{G})$
we choose an open covering $\mathcal{U}$ and an element
$\overline{h} \in \check{H}^n(\mathcal{U}, \mathcal{G})$
mapping to $\xi$ in the colimit defining {\v C}ech cohomology.
Then we choose $\mathcal{V}$ and $g$ as above and set
$\delta(\xi)$ equal to the image of $\delta(\overline{h})$
in $\check{H}^n(X, \mathcal{F})$.
At this point a lot of properties have to be checked, all of which
are straightforward. For example, we need to check that our construction
is independent of the choice of
$\mathcal{U}, \overline{h}, \mathcal{V}, \alpha : J \to I, g$.
The class of $\text{d}(g)$ is independent of the choice of the lifts
$g_{i_0 \ldots i_n}$ because the difference will be a coboundary.
Independence of $\alpha$ holds\footnote{This is an important
check because the nonuniqueness of $\alpha$ is the only thing preventing
us from taking the colimit of {\v C}ech complexes over all open
coverings of $X$ to get a short exact sequence of complexes computing
{\v C}ech cohomology.}
because a different choice
of $\alpha$ determines homotopic vertical maps of complexes
in the diagram above, see Section \ref{section-refinements-cech}.
For the other choices we use that given a finite collection
of open coverings of $X$ we can always find an open covering refining all
of them. We also need to check additivity which is shown in the same manner.
Finally, we need to check that the maps
$\check{H}^n(X, -) \to H^n(X, -)$ are compatible
with boundary maps. To do this we choose injective
resolutions
$$
\xymatrix{
0 \ar[r] &
\mathcal{F} \ar[r] \ar[d] &
\mathcal{G} \ar[r] \ar[d] &
\mathcal{H} \ar[r] \ar[d] &
0 \\
0 \ar[r] &
\mathcal{I}_1^\bullet \ar[r] &
\mathcal{I}_2^\bullet \ar[r] &
\mathcal{I}_3^\bullet \ar[r] &
0
}
$$
as in Derived Categories, Lemma \ref{derived-lemma-injective-resolution-ses}.
This will give a commutative diagram
$$
\xymatrix{
0 \ar[r] &
\check{\mathcal{C}}^\bullet(\mathcal{U}, \mathcal{F}) \ar[r] \ar[d] &
\check{\mathcal{C}}^\bullet(\mathcal{U}, \mathcal{F}) \ar[r] \ar[d] &
\check{\mathcal{C}}^\bullet(\mathcal{U}, \mathcal{F}) \ar[r] \ar[d] &
0 \\
0 \ar[r] &
\text{Tot}(\check{\mathcal{C}}^\bullet(\mathcal{U}, \mathcal{I}_1^\bullet))
\ar[r] &
\text{Tot}(\check{\mathcal{C}}^\bullet(\mathcal{U}, \mathcal{I}_2^\bullet))
\ar[r] &
\text{Tot}(\check{\mathcal{C}}^\bullet(\mathcal{U}, \mathcal{I}_3^\bullet))
\ar[r] &
0
}
$$
Here $\mathcal{U}$ is an open covering as above and
the vertical maps are those used to define the maps
$\check{H}^n(\mathcal{U}, -) \to H^n(X, -)$, see
Lemma \ref{lemma-cech-cohomology}.
The bottom complex is exact as the sequence of
complexes of injectives is termwise split exact.
Hence the boundary map in cohomology is computed
by the usual procedure for this lower exact sequence, see
Homology, Lemma \ref{homology-lemma-long-exact-sequence-cochain}.
The same will be true after passing to the refinement
$\mathcal{V}$ where the boundary map for {\v C}ech cohomology
was defined. Hence the boundary maps agree because they
use the same construction (whenever the first one is defined
on an element in {\v C}ech cohomology on a given open covering).
This finishes our discussion of the construction of
the structure of a $\delta$-functor on {\v C}ech cohomology
and why this structure is compatible with the given
$\delta$-functor structure on usual cohomology.

\medskip\noindent
Finally, we may apply Lemma \ref{lemma-injective-trivial-cech}
to see that higher {\v C}ech cohomology is trivial on injective
sheaves. Hence we see that {\v C}ech cohomology is a universal
$\delta$-functor by
Homology, Lemma \ref{homology-lemma-efface-implies-universal}.
\end{proof}
